\documentclass[11pt]{article}
\usepackage[utf8]{inputenc}
\usepackage[T1]{fontenc}
\usepackage{lmodern}
\usepackage[margin=1in]{geometry}
\usepackage{amsmath,amssymb}
\usepackage{textcomp}
\usepackage[dvipsnames]{xcolor}
\usepackage{enumitem}
\usepackage{parskip}
\usepackage{titlesec}
\usepackage[colorlinks=true,linkcolor=blue,citecolor=blue,urlcolor=blue]{hyperref}
\usepackage{natbib}

\titleformat{\section}{\large\bfseries}{\thesection}{1em}{}
\titleformat{\subsection}{\normalsize\bfseries}{\thesubsection}{1em}{}
\titleformat{\subsubsection}{\normalsize\bfseries\itshape}{}{0em}{}

\title{\textbf{Not All Domains Degrade Equally:\\ Correctness Strictness and Synthetic-Contamination Fragility}}
\author{}
\date{}

\newcommand{\ph}{\textit{[placeholder --- fill in after experiments are run]}}

\begin{document}
\maketitle
\vspace{-2em}

\begin{quote}
\small\textit{Assembled from the six confirmed stage outputs of an \texttt{auto-research-discussion} run (seed topic B6.4), post research-taste audit (\textcolor{ForestGreen}{\textbf{PASS}}), post stage review (1 PASS, 3 REVISE-fixed, 1 PASS). Status: \textbf{specification / paper draft --- no experiments have been run.} All numbers below are pre-registered design targets, not results.}
\end{quote}

\begin{abstract}
Large language models increasingly pretrain on web corpora already mixed with model-generated text, and practitioners with a finite budget for detecting or filtering that content need to know where it hurts most. We show this question cannot currently be answered by either of the two literatures that come closest: model-collapse theory expresses contamination as a single global mixing ratio with no domain axis, while data-mixture optimization treats per-domain composition as a causally live lever but assumes every domain shard is clean. We hypothesize instead that damage concentrates by a domain's correctness structure: domains with mechanically verifiable ground truth (code, mathematics) may convert a subtly wrong synthetic document into corrupted training signal, while style-tolerant domains (prose, dialogue) may absorb the same dose as benign variation. We build a token-accounted, single-domain contamination protocol that normalizes capability damage across domains with incommensurable metrics, isolates syntheticness from co-occurring textual degeneracy via a matched real-text control, and rules out generator-specific artifacts via cross-generator replication. We report whether contamination damage is domain-selective, which domains are most fragile, and whether the resulting ranking survives scale, generator, and degeneracy controls --- evidence intended to directly guide synthetic-data-filtering priority in practice.
\end{abstract}

\section{Introduction}

Synthetic text generated by large language models is no longer a rare visitor to web-scale pretraining corpora --- it arrives unlabeled, at scale, and increasingly by default, as more of the web is itself written with LLM assistance. Teams doing continued pretraining or domain-adaptive fine-tuning on such corpora face a concrete, recurring decision: with a finite budget for detecting or filtering synthetic content, where should that effort go first? The intuitive answer --- treat contamination as a single global quality dial and filter uniformly --- is cheap to implement, which is exactly why it is worth checking rather than assuming.

Yet a specific question about this contamination remains unaddressed: does the same dose of synthetic text injected into different parts of a training corpus cause the same amount of downstream damage, or does harm concentrate in particular content domains? Put plainly, damage from synthetic contamination may not be a property of the corpus as a whole but of the domain it lands in --- and this is not a rhetorical question. A domain whose correctness can be mechanically verified --- code that must execute, mathematics whose final answer can be checked --- may have no defense against a synthetic document that is subtly wrong: the error passes as plausible text and becomes corrupted supervision the model has no signal to discount. A domain whose quality is judged rather than checked --- prose, casual dialogue --- may absorb the same synthetic dose as ordinary stylistic variation. If this asymmetry is real, a domain-uniform filtering policy --- the only kind current theory can even express --- systematically over-protects the domains that can tolerate contamination and under-protects the ones that cannot.

The nearest existing work approaches this gap from two adjacent but non-overlapping angles, and it is worth naming the closest pair plainly before surveying the wider landscape: work on preventing model collapse, such as the mixing-ratio analysis of \citet{MixRatio2023}, formalizes contamination risk as a function of a single global real-synthetic mixing weight over one homogeneous distribution --- its own framework has no axis along which a domain could even be named, let alone compared to another. Data-mixture optimization, such as \citet{DoReMi2023}, sits on the other side of the same gap: it proves that per-domain composition materially shifts downstream capability, establishing that the domain axis is causally live, but every domain shard it reweights is assumed clean. Between them, these two works bound the exact question this paper asks --- contamination varies, and domain varies --- but neither work, nor their union, crosses the two axes against each other. The closest empirical datum to date only sharpens the gap rather than closing it: an epidemic-style model of contamination dynamics \citep{EpidemiologyCollapse2026} observes, as a robustness aside, that its dose-response pattern replicates across the two small corpora it tests --- but under a homogeneous-mixing framework, at GPT-2 scale, without decomposing a mixed corpus by content domain, so whether that apparent uniformity survives a domain-decomposed design at capability-relevant scale is exactly what remains unmeasured.

Turning this asymmetry into a measurable claim surfaces its own obstacles. The four candidate domains admit no common unit of damage: code's executable pass/fail criterion, mathematics's checkable final answer, and prose or dialogue's graded or judged scores cannot be compared as raw numbers without the comparison itself silently determining which domain looks most fragile. Within any single contaminated cell, a synthetic document's ``syntheticness'' and its incidental textual degeneracy --- repetition, reduced lexical variety --- change together, so any observed damage is ambiguous between the two causes unless a control is built to separate them. Across cells, the picture is further complicated by scale: the full grid of domain, contamination ratio, and model size is combinatorially expensive at pretraining-relevant compute, so any affordable study samples a small corner of that space and must justify why the untested corners would not overturn its conclusion. And any ranking built from one contamination generator is entangled with that generator's own domain-dependent writing quality --- a generator that happens to write weaker code than prose would make code look artificially fragile regardless of any domain-intrinsic effect.

Our approach turns each of these obstacles into a specific design choice rather than a hedge. Every domain's capability change is scored as a delta normalized against its own clean-control baseline, computed independently under both a domain-native benchmark and a benchmark-independent validation-loss measure, so cross-domain comparison is possible without one metric's sensitivity silently deciding the ranking. Syntheticness is isolated from incidental degeneracy by comparing the contaminated arm against a real-text control built to match the synthetic arm's measured repetition and lexical-variety profile, so any excess damage beyond that match is not explained by the two degeneracy covariates specifically. A staged design tests dose-response shape at full resolution at one model scale, then replicates only the two most informative ratio extremes at a second, larger scale, checking whether the resulting domain ranking --- not its exact magnitude --- transfers across scale. And the entire sweep is independently repeated with a second contamination generator from a different model family, so a ranking that survives the swap cannot be explained by either generator's idiosyncratic per-domain skill.

We evaluate this design across four domains --- code, mathematics, prose, and dialogue --- drawn from public per-domain corpora, contaminating one domain shard at a time against a shared clean control, at model scales up to the low single-digit billions of parameters.

In summary, we contribute:
\begin{itemize}[leftmargin=1.4em]
\item An identification and formalization of a domain-selectivity hypothesis for synthetic contamination: damage may track a domain's correctness strictness rather than acting as the domain-uniform quality tax current model-collapse theory implicitly assumes.
\item A dual-normalization measurement instrument that makes capability damage comparable across domains with fundamentally incommensurable ground-truth structures, separating the claim that damage is domain-selective at all from the claim that it is ordered by correctness strictness specifically.
\item A degeneracy-matched isolation test and a cross-generator replication test that together bound how much of any observed ranking is attributable to syntheticness itself rather than to incidental text quality or to one generator's idiosyncrasies --- reported as a partial, not total, identification.
\item A staged dose-response and cross-scale design that tests whether a domain vulnerability ranking, rather than its exact magnitude, transfers from small-scale proxy experiments toward the scale at which curation decisions actually get made.
\end{itemize}

\section{Problem Definition}

\noindent\textbf{Formal setup.} Let the pretraining corpus be $D = \bigcup_{d \in \mathbb{D}} D_d$, with $\mathbb{D} = \{\text{code}, \text{math}, \text{prose}, \text{dialogue}\}$ denoting four domain shards. A contamination generator $G$ produces in-domain synthetic documents. The contamination operator $\kappa(D, d, r)$ replaces exactly an $r$-fraction of shard $D_d$'s tokens with $G$-synthetic documents, holding every other shard byte-identical and the total token count fixed. A model $\theta(d, r, s)$ is continued-pretrained on $\kappa(D, d, r)$ with random seed $s$ under an otherwise identical schedule. $S_d(\cdot)$ denotes a domain-$d$ capability metric evaluated on a checkpoint. The normalized capability delta is $\Delta_d(r)$, the (sign-adjusted) fractional loss of $S_d(\theta(d,r,s))$ relative to $S_d(\theta(d,0,s))$, and $\beta_d$ is the ordinary-least-squares slope of $\Delta_d$ on $r$ across ratios and seeds. The object of study is whether the domain $\times$ ratio interaction on $\{\Delta_d(r)\}$ is significant, and if so, whether the resulting ordering of $\{\beta_d\}$ is directionally consistent with correctness strictness, generator-stable, and robust to a degeneracy-matched control.

\noindent\textbf{Why this is hard.} The four domains in $\mathbb{D}$ admit no common raw unit --- $S_{\text{code}}$ is an execution pass rate, $S_{\text{math}}$ is exact-match accuracy, $S_{\text{prose}}$ and $S_{\text{dialogue}}$ are perplexity or judged scores --- so comparing raw values of $S_d$ across $d$ would let an accident of benchmark sensitivity determine the apparent ranking; a valid comparison requires normalizing every $S_d$ against its own $\theta(d,0,s)$ baseline, and even that normalization choice is itself a modeling decision whose effect on the ranking must be checked, not assumed away. Within a single cell $\kappa(D,d,r)$, the injected documents' syntheticness and their incidental textual degeneracy vary together by construction, so $\Delta_d(r)$ alone cannot attribute damage to either cause specifically. Across cells, the joint space of $(d, r, \text{model scale})$ grows multiplicatively, so any affordable design samples only part of it and must justify extrapolation from the sampled corner to the scale at which curation decisions are actually made. And because $\kappa(D,d,r)$ depends on a specific choice of $G$, any ordering of $\{\beta_d\}$ obtained from one generator is confounded with that generator's own domain-dependent writing quality unless the ordering is re-derived under an independently-sourced second generator.

\section{Method Overview}

At a conceptual level, the study is a single measurement pipeline with five cooperating parts, each answering one of the obstacles above. Domain shards and a contamination generator enter the pipeline; a dual-normalization scorer turns the resulting per-domain checkpoints into comparable capability deltas and a fitted interaction test; a degeneracy-matched control arm turns the same deltas into a syntheticness-specific damage estimate; a staged ratio-and-scale design turns the deltas into a dose-response characterization whose ranking is checked for scale-stability; a second, independently-sourced generator turns the whole exercise into a claim about domain identity rather than one generator's idiosyncrasy; and a verifiable-error audit of the injected documents themselves turns the aggregate ranking into an independently-checkable claim about \emph{why} strict-correctness domains are fragile, not only \emph{that} they are.

The \textbf{dual-normalization capability scorer} takes each domain's raw benchmark score and its held-out validation loss and produces two independent normalized deltas per checkpoint, fitting a single mixed-effects model whose interaction term answers whether damage is domain-selective at all and whose fitted per-domain slopes answer, separately, whether the ordering follows correctness strictness. The \textbf{degeneracy-matched syntheticness isolator} takes the contaminated arm and builds a matched real-text arm with the same measured repetition and lexical-variety profile, attributing any damage beyond that match to syntheticness rather than incidental quality loss --- while explicitly not claiming to rule out unmeasured, latent quality defects. The \textbf{staged dose-response and scale design} takes a coarse two-point sweep at one model scale to establish that an effect exists, a full ratio grid at that same scale to characterize its shape, and a replication at two ratio extremes at a larger scale to test whether the resulting ranking, not its magnitude, transfers. The \textbf{cross-generator replication test} takes the entire sweep and repeats it with a second, structurally different generator, comparing the two resulting rankings for agreement. The \textbf{verifiable-error audit} takes the synthetic documents actually injected into the code and mathematics shards, checks each one mechanically (execution, symbolic answer-checking) for correctness independent of surface fluency, and correlates the resulting per-domain error rate against the fitted slopes --- explanatory evidence for the mechanism, reported separately from the ranking's own pass/fail criteria.

\section{Methodology}

\subsection{Dual-Normalization Capability Scorer}
\textbf{Challenge.} The four domains have no common raw unit of capability, so comparing benchmark scores directly across domains would let benchmark sensitivity, not domain fragility, determine the apparent ranking.

\textbf{Insight.} Normalizing every domain against its own clean-control baseline makes cross-domain comparison possible in the first place, but the direction of any resulting ranking is a separate, substantive claim the normalization does not itself guarantee --- one measurement apparatus yields two logically distinct, separately falsifiable claims.

\textbf{Method.} Each domain's capability change is expressed as a delta normalized against its own zero-contamination baseline, computed independently under the domain-native benchmark metric and under a benchmark-independent held-out validation-loss metric. A mixed-effects model regresses the normalized delta on contamination ratio interacted with domain, with seed as a random effect. The interaction term tests whether damage is domain-selective at all; a pre-registered contrast between the mean slope of code and mathematics versus the mean slope of prose and dialogue, checked under both normalizations, tests the stronger, directional correctness-strictness hypothesis.

\subsection{Degeneracy-Matched Syntheticness Isolation}
\textbf{Challenge.} Within a single contaminated cell, a synthetic document's syntheticness and its incidental textual degeneracy (repetition, reduced lexical variety) change together, so observed damage cannot be attributed to either cause on its own.

\textbf{Insight.} Degeneracy and syntheticness are inseparable inside one cell, but they are separable across two cells deliberately built to hold measured degeneracy fixed while syntheticness differs --- though only for the degeneracy that is actually measured, not for every possible latent quality defect.

\textbf{Method.} Per-document degeneracy covariates (generator perplexity, type-token ratio) are logged during contamination injection. A matched control arm replaces the same token budget with real, non-synthetic documents from the same domain, subsampled or filtered to match the synthetic arm's measured repetition and type-token-ratio profile. Damage in the synthetic arm is compared against this degeneracy-matched real-text arm; any excess is attributed to syntheticness beyond the two logged covariates specifically. This is reported as a \textbf{partial, not total, identification}: the comparison equates measured surface degeneracy between the two arms, but it cannot equate latent degeneracies the logged covariates do not capture --- for example, subtle factual or logical drift --- so some residual confounding between syntheticness and unmeasured quality defects remains even after the comparison.

\subsection{Staged Dose-Response and Scale-Stability Design}
\textbf{Challenge.} The joint space of domain, contamination ratio, and model scale grows multiplicatively, and exhaustively training it is unaffordable at pretraining-relevant compute.

\textbf{Insight.} The paper's actionable claim only requires the domain \emph{ranking}, not the exact damage \emph{magnitude}, to transfer across scale --- so scarce large-scale compute can be concentrated on the two most informative ratio extremes rather than spread thinly across a full grid, while a full ratio sweep at a smaller, affordable scale still characterizes dose-response shape.

\textbf{Method.} A coarse two-point ratio sweep at pilot scale establishes whether an effect exists and sizes it. A full five-point ratio grid at base model scale characterizes each domain's dose-response shape. The two most extreme ratios are then replicated at a second, larger model scale, and the resulting per-domain slope ranking --- not the raw slope values --- is compared across the two scales via Kendall's tau.

\subsection{Cross-Generator Replication}
\textbf{Challenge.} Any domain vulnerability ranking obtained from one contamination generator is entangled with that generator's own domain-dependent writing quality, and no single-generator study can separate the two by construction.

\textbf{Insight.} A generator-specific artifact has no reason to reproduce under a structurally different generator, so requiring two independently-sourced generators to agree on the ranking is the only available test that can rule out ``this ranking is really a fingerprint of generator X.''

\textbf{Method.} The entire single-domain contamination sweep is independently repeated using a second contamination generator from a different model family and quality tier, holding the injection protocol (in-domain conditioning, near-duplicate rejection, length matching) fixed. Agreement between the two generators' resulting per-domain slope rankings is measured via Kendall's tau.

\subsection{Mechanism Analysis: A Verifiable-Error Audit}
\textbf{Challenge.} The measurement apparatus above establishes \emph{that} damage is domain-selective and, if so, \emph{which} direction it runs --- but neither test probes \emph{why} a strict-correctness domain would be more fragile in the first place; without an independent mechanism check, the correctness-strictness story remains an untested account bolted onto an aggregate ranking.

\textbf{Insight.} The Step-0 mechanism claim is that a domain's mechanically verifiable correctness structure lets subtly wrong synthetic content pass as plausible text and become corrupted supervision. That claim makes an independent, checkable prediction distinct from the ranking itself: the rate at which injected synthetic documents are themselves verifiably wrong should track how much damage that domain sustains, not merely correlate with the ranking by construction.

\textbf{Method.} For code, each injected synthetic document's associated unit tests (where authored) or a reference solution's tests are executed, and the fraction that is verifiably incorrect despite passing surface fluency is recorded. For mathematics, each synthetic worked solution's final answer is resolved against a symbolic or numeric checker and the fraction with a wrong answer is recorded. Prose and dialogue admit no analogous mechanical checker by construction, so their audit is fixed at ``not computable'' rather than assumed zero. The resulting per-domain verifiable-error rate (where computable) is correlated against $\beta_d$ across seeds and ratios (Spearman $\rho$), reported descriptively as explanatory evidence for \emph{why} RQ1/RQ2 hold rather than as a third pass/fail gate: a null or negative correlation would not overturn RQ1/RQ2 if they pass on their own pre-registered criteria, but would flag that the correctness-strictness \emph{mechanism} is unsupported even where the \emph{ranking} holds --- the descriptive (ranking) and causal (mechanism) claims are recorded as separable, not conflated.

\section{Experiments}

\subsection{Experimental Setup}

\textbf{Feasibility check.} Before committing to the full design, a three-arm walking skeleton (clean control; code at full contamination; prose at full contamination; one seed each, roughly a day on a four-GPU node) establishes that at least one domain shows a measurable, domain-localized delta at full contamination through the entire pipeline, using only deterministic metrics (execution-based pass rate, last-word accuracy, held-out validation loss --- no LLM judge). Scaling to the full pilot and to the scale-up grid changes only the configuration matrix (ratios, seeds, generator identity, model size), not the code path; if the skeleton shows no domain-localized effect at full contamination, the full grid is not run.

\textbf{Compute budget} (pre-registered design targets, not measurements; per-run cost anchored to the walking skeleton's estimate of 6--8 A100-hours per 1.1B-parameter / 0.5B-token continued-pretraining run). Walking skeleton: 3 runs $\approx$ 20--25 A100-hours. Pilot: 27 runs (4 domains $\times$ $r \in \{25\%, 100\%\}$ $\times$ 3 seeds, plus 3 shared clean controls). Stage-2 additions at base scale: the two remaining ratio points (24 runs), the second-generator full sweep (48 runs), the degeneracy-matched real-text arm at full contamination (12 runs), and the global-uniform arm (3 runs) --- $\approx$117 continued-pretraining runs at 1.1B scale in total, $\approx$700--950 A100-hours. Scale-up: the 3B replication at the two ratio extremes adds 15 runs per generator at roughly three times the per-run cost ($\approx$300--600 A100-hours depending on whether the second generator is also replicated at scale). End to end, including synthetic-data generation and evaluation overhead, the full design budgets $\approx$1,000--1,500 A100-hours --- which is why RQ3's scale leg is restricted to the two most informative ratio extremes and RQ4 reuses the base scale only, and what ``at the edge of this study's compute budget'' means concretely in the Further-evaluation note below.

\textbf{Datasets.} The Stack (permissive-license Python/Java subset) --- code training shard. OpenWebMath --- mathematics training shard. FineWeb-Edu --- prose training shard. OASST1 combined with UltraChat --- dialogue training shard, used as the ``clean'' baseline \textbf{contingent on a pre-registered audit}: a synthetic-text classifier is passed over the pooled shard before any contamination run to estimate its existing synthetic base rate, since UltraChat is itself LLM-generated and an undetected synthetic base rate would silently compress the measurable dialogue-domain delta. If the audit finds a base rate above 20\%, UltraChat is dropped in favor of OASST1 plus a human-only forum corpus as a pre-committed design amendment, before the dialogue domain enters the ratio sweep at all. The code shard's evidence table (paired technique comparisons, ablation-style rows) follows the structural pattern of \citet{EffectivenessTrainingDat}'s controlled empirical study of code-domain training-data-optimization techniques, imitated for shape (not content) per this paper's own structural-exemplar discipline.

\textbf{Baselines} (sanity-check comparators, not competitors to beat --- this is a measurement study, and no research question is framed as ``our procedure beats baseline X''). A clean-control arm (0\% contamination, shared across domains per seed) is the necessary reference for every contaminated cell. A global-uniform contamination arm --- the same total synthetic-token budget spread proportionally across all four shards --- is the ``why do this at all'' control: if concentrated single-domain injection is statistically indistinguishable from proportionally-spread injection, the domain axis adds nothing and the problem this paper studies does not exist. The degeneracy-matched real-text arm and the second-generator arm are used directly by the methods above. Two optional mitigation-characterization arms --- an entropy-maximizing selection heuristic in the style of \citet{EntropyCollapse2023}, and an importance-weighted synthetic-training heuristic in the style of \citet{SynthWeighting2023} --- characterize how much of any measured damage an existing generic mitigation already recovers, again strictly as sanity checks rather than competitors.

\textbf{Metrics.} Domain-native capability scores (execution-based pass rate for code, exact-match accuracy for mathematics, last-word accuracy and perplexity for prose, judged score for dialogue), each domain's normalized capability delta under both the benchmark and validation-loss normalizations, per-domain OLS slopes with bootstrap confidence intervals, and Kendall's tau between rankings under different conditions (scale, generator). The dialogue judged score is the design's only non-deterministic metric, and it is gated by a pre-registered validity check before any use at scale: a frozen single-turn 1--10 judge is validated on 100 dialogue responses spanning clean and contaminated arms, each double-rated by two human raters on the same rubric, requiring Cohen's $\kappa \geq 0.7$ between the judge and the human-rater majority, computed once and frozen before the ratio sweep. If the threshold is not met, the dialogue domain's primary metric demotes to held-out validation loss only (already co-primary), with the judge score reported descriptively and excluded from the Holm-corrected primary-metric family --- dialogue leaves the judged-metric family, not the study.

\textbf{Research questions:}
\begin{enumerate}[leftmargin=1.6em]
\item RQ1 --- Is contamination damage domain-selective at all?
\item RQ2 --- Is the domain ordering consistent with correctness strictness (code and mathematics more fragile than prose and dialogue)?
\item RQ3 --- Does dose-response scale sensibly with contamination ratio, and does the resulting ranking transfer across model scale?
\item RQ4 --- Does the ranking replicate under an independently-sourced second contamination generator?
\item RQ5 --- Does damage in the synthetic arm exceed a degeneracy-matched real-text control?
\end{enumerate}

\subsubsection{RQ1: Is contamination damage domain-selective at all?}
\textbf{Setup.} Fit the dual-normalized mixed-effects model (normalized delta $\sim$ domain $\times$ ratio $+ (1 \mid \text{seed})$) on the pilot and scale-up grids; test the domain $\times$ ratio interaction term via a likelihood-ratio test against a no-interaction model, Holm-corrected across the four domain-matched primary metrics, $\alpha = 0.05$.

\textbf{Expected outcome.} The interaction term is significant, indicating that damage is not a domain-uniform tax.

\textbf{Falsification condition.} Refuted if the interaction term is non-significant after Holm correction. In that case, the reportable finding reframes to: within the tested regime (up to the low single-digit billions of parameters, up to 100\% single-shard contamination), synthetic contamination behaves as a domain-uniform capability tax, with every cross-domain slope difference bounded below the observed upper confidence limit --- a result that would simplify, rather than complicate, existing global-mixing-ratio theory.

\textbf{Hypothesis (plain language).} Contamination damage concentrates unevenly across domains rather than acting as a uniform tax on the whole corpus.

\textbf{Metrics.} Interaction-term likelihood-ratio statistic and Holm-corrected $p$-value; per-domain slopes and their pairwise confidence intervals.

\textbf{Result.} \ph

\subsubsection{RQ2: Is the domain ordering consistent with correctness strictness?}
\textbf{Setup.} Conditional on RQ1, test the pre-registered contrast (mean slope of code and mathematics) minus (mean slope of prose and dialogue) via a one-sided bootstrap 95\% confidence interval, checked independently under the benchmark-metric and validation-loss normalizations.

\textbf{Expected outcome.} The contrast is positive and its confidence interval excludes zero under both normalizations, indicating that domains with mechanically verifiable correctness structure are more fragile to contamination.

\textbf{Falsification condition.} Refuted if the contrast's confidence interval includes zero under either normalization, or if the sign is negative --- either would mean the correctness-strictness account of \emph{which} domains are more fragile is wrong, even if RQ1's domain-selectivity finding still holds. In that case, the paper reports the observed ranking with its confidence intervals as an open empirical finding without the strictness explanation, noting that any stable ranking still supports a curation-priority recommendation, just not the correctness-strictness mechanism story.

\textbf{Hypothesis (plain language).} Domains whose ground truth can be mechanically checked are more vulnerable to subtly wrong synthetic content than domains whose quality is judged rather than checked.

\textbf{Metrics.} Sign and 95\% bootstrap confidence interval of the strict-versus-tolerant slope contrast under both normalizations.

\textbf{Result.} \ph

\subsubsection{RQ3: Does dose-response scale sensibly, and does the ranking transfer across model scale?}
\textbf{Setup.} Fit per-domain OLS slopes and quadratic terms across the full ratio grid at base scale; replicate the two most extreme ratios at a second, larger model scale; compare the domain rankings at the two scales via Kendall's tau.

\textbf{Expected outcome.} The normalized delta is non-decreasing in contamination ratio in at least three of four domains, and the domain ranking agrees across the two tested scales at Kendall's tau of at least 0.66 (at most one adjacent transposition among four domains).

\textbf{Falsification condition.} Refuted if dose-response is flat or non-monotonic in a way that breaks the ranking's ordering, or if Kendall's tau between the two scales falls below 0.66. In that case, the vulnerability ranking is reported as \textbf{scale-conditional} --- it holds at the smaller tested scale but reorders at the larger one --- meaning curation priorities derived from small-scale proxy studies would not transfer to the scale at which curation decisions are actually made.

\textbf{Hypothesis (plain language).} The domain ranking is not an artifact of the specific small model scale used to produce it.

\textbf{Metrics.} Per-domain slope and quadratic-term estimates; cross-scale Kendall's tau on the extreme-ratio ranking.

\textbf{Result.} \ph

\subsubsection{RQ4: Does the ranking replicate under a second, independently-sourced contamination generator?}
\textbf{Setup.} Repeat the full single-domain contamination sweep with a second generator from a different model family and quality tier, holding the injection protocol fixed; compute Kendall's tau between the two generators' resulting per-domain slope rankings.

\textbf{Expected outcome.} Kendall's tau between the two generators' rankings is at least 0.66.

\textbf{Falsification condition.} Refuted if Kendall's tau falls below 0.66. In that case, per-domain contamination vulnerability is reported as \textbf{generator-conditional} --- rankings from any single generator are unreliable on their own, which would make cross-generator replication a mandatory step for any future contamination audit rather than an optional robustness check, and the measurement protocol itself (rather than any specific ranking) becomes the paper's most defensible contribution.

\textbf{Hypothesis (plain language).} The observed ranking is a property of the domains, not a fingerprint of the specific generator used to contaminate them.

\textbf{Metrics.} Kendall's tau between generator-wise rankings, with bootstrap confidence interval over seeds and evaluation items.

\textbf{Result.} \ph

\subsubsection{RQ5: Does damage in the synthetic arm exceed a degeneracy-matched real-text control?}
\textbf{Setup.} At full contamination, compare the synthetic arm's capability damage against the degeneracy-matched real-text arm via a paired one-sided bootstrap test per domain, Holm-corrected across the four domains.

\textbf{Expected outcome.} The synthetic arm shows significantly more damage than the degeneracy-matched arm in at least half of the tested domains, including the single most vulnerable domain identified by the cross-domain ranking.

\textbf{Falsification condition.} Refuted if the degeneracy-matched arm degrades as much as the synthetic arm across the tested domains. In that case, contamination damage in this regime is explained by \textbf{surface degeneracy rather than syntheticness specifically} --- a budget-relevant result in its own right, since it would mean cheap degeneracy detectors (repetition or lexical-variety filters) recover most of the protection that expensive synthetic-text classifiers would otherwise provide.

\textbf{Hypothesis (plain language).} At least part of the measured damage is attributable to a document being machine-generated, not only to the generic textual degeneracy such documents happen to carry.

\textbf{Metrics.} Paired one-sided bootstrap $p$-values per domain (Holm-corrected); effect size of the synthetic-versus-degeneracy-matched contrast.

\textbf{Result.} \ph

\subsubsection{Mechanism check: does verifiable error track the ranking?}
\textbf{Setup.} For the code and mathematics shards, execute each injected synthetic document's associated tests or a symbolic/numeric answer checker, and record the fraction verifiably incorrect despite passing surface fluency; correlate this per-domain error rate against the fitted slopes $\beta_d$ across seeds and ratios via Spearman $\rho$. Prose and dialogue have no analogous mechanical checker and are recorded as ``not computable'' rather than assumed null.

\textbf{Status.} This is explanatory, not a sixth pass/fail gate: it is reported descriptively alongside RQ1/RQ2 rather than aggregated into their Holm correction. A null or negative correlation would not overturn RQ1/RQ2 if those pass on their own pre-registered criteria, but would flag that the correctness-strictness \emph{mechanism} account is unsupported even where the \emph{ranking} holds --- the two are recorded as separable findings.

\textbf{Result.} \ph

\vspace{0.5em}
\noindent\textbf{Further evaluation.} With more time, a natural extension is testing generalization beyond the four tested domains (e.g., legal or biomedical text, whose correctness structure sits between the tested extremes) and beyond the tested model-scale ceiling --- both explicitly out of scope here because the current design's scale-stability check (RQ3) and generator-stability check (RQ4) are themselves already at the edge of this study's compute budget.

\section{Related Work}

The two nearest prior works are MixRatio \citep{MixRatio2023} and DoReMi \citep{DoReMi2023}, and it is worth stating the delta plainly before surveying the wider landscape: MixRatio already formalizes contamination risk as a function of a real-synthetic mixing weight, and DoReMi already proves that per-domain composition materially shifts downstream capability --- but MixRatio's framework has no domain axis to decompose, and DoReMi's domain axis is assumed clean rather than contaminated. Together they bound exactly the gap this paper closes; neither, nor their union, crosses the contamination axis against the domain axis.

\textbf{Model-collapse and recursive-training theory.} \citet{MixRatio2023} derives closed-form long-run risk as a function of a real-synthetic mixing weight in an overparameterized regime; \citet{IterStability2023} derives contraction conditions under which iterative retraining on mixed real/synthetic data remains stable rather than collapsing; \citet{CrossModelRecursion2023} models multi-entity recursive training through a shared synthetic-data channel and documents aggregate tail-event loss; \citet{EntropyCollapse2023} reinterprets collapse as a generalization-to-memorization transition driven by declining dataset entropy. \citet{Shumailov2024Collapse} established the foundational empirical and theoretical demonstration that recursive training on model-generated content causes irreversible distributional defects. All five model contamination as a property of one homogeneous distribution or generation-indexed pool; none decomposes the corpus into content domains, so none can express --- let alone rank --- differential per-domain fragility.

\textbf{Data-mixture optimization.} \citet{DoReMi2023} trains a small proxy model with a Group-DRO objective over domain excess losses to derive mixture weights that transfer to large-scale pretraining; \citet{RegMix2023} trains many small proxy models on sampled domain mixtures and regresses validation performance on mixture proportions, exploiting rank-invariance of mixture quality across scales --- a rank-transfer logic this paper's own cross-scale design directly borrows. Both establish that per-domain composition materially shifts downstream capability, proving the domain axis is causally live, but both explicitly assume the domain shards being weighted are clean; neither includes a contamination axis, so neither says anything about where synthetic contamination hurts most. \citet{CAUSALM} independently reframes domain-mixture selection as a causal-inference problem and \citet{CanSmallTraining} interrogates whether small-scale proxy-model mixture rankings reliably transfer to full-scale training --- a validity question this paper's own cross-scale replication (RQ3) directly inherits and tests in the contaminated setting rather than the clean-mixture setting both papers study.

\textbf{Synthetic-data reweighting and domain-specific fine-tuning.} \citet{SynthWeighting2023} reweights synthetic fine-tuning examples for one task and domain at a time via closed-form importance weights correcting distribution mismatch against a small real dataset; \citet{TALR2023} studies how domain-specific supervised fine-tuning on real, curated data perturbs general capability through optimization dynamics, showing small learning rates and token-adaptive reweighting preserve general capability. \citet{SpendWisely2023} separately optimizes synthetic-sample-count schedules across bootstrapping iterations for one task-homogeneous data pool. \citet{SynthWeighting2023} operates within a single task and domain at fine-tuning scale rather than comparing fragility across domains; \citet{TALR2023}'s training data is real rather than synthetic, so it cannot separate a domain's own training effect from a synthetic contamination effect; and \citet{SpendWisely2023} offers no guidance on staging a domain-decomposed grid under a fixed compute budget. None of the three provides evidence about which domain is most vulnerable to synthetic contamination or why.

\textbf{Concurrent and subsequent work (as of 2026-08).} Three lines of 2025--26 work now sit adjacent to this paper's claims and are conceded plainly. First, an epidemic-style model of contamination spread \citep{EpidemiologyCollapse2026} couples corpora and models as interacting populations and --- in a robustness aside --- reports that its measured dose-response pattern replicates across the two small corpora it tests; this is the closest existing empirical datum on RQ1, and it points toward the domain-uniform null, but under homogeneous mixing, at GPT-2 scale, with two arbitrary text corpora rather than a content-domain decomposition --- precisely the design this paper pre-registers. Second, a study of recursive self-training collapse in code LLMs \citep{CodeRecursiveCollapse2026} documents collapse dynamics and review-gate failure modes in the code domain alone; it occupies the code-only corner of this paper's grid and is consistent with the correctness-strictness intuition, but with no second domain it cannot say whether code is \emph{differentially} fragile --- the comparative estimand is untouched. Third, the largest-scale factorial study of synthetic data in pretraining to date (arXiv:2510.01631; over 1,000 models) crosses synthetic-data \emph{type} (rephrased vs.\ textbook-style) with mixing ratio and finds collapse evidence is type-dependent; its manipulated axis is how the synthetic text was produced, not which content domain it lands in, and it never injects into a single domain shard. A further verification-based escape result \citep{VerifierEscape2025} shows an external verifier prevents collapse but bounds gains at the verifier's own level, joining \citet{IterativeDeploymentImpro} as a second contrast case on the strict-verifiability axis this paper's hypothesis is built on. None of the four decomposes a mixed corpus by target content domain, injects into one shard at a time, or estimates a cross-domain fragility ranking.

\textbf{A contrast case for the correctness-strictness hypothesis.} \citet{IterativeDeploymentImpro} reports the opposite outcome from the one this paper's central hypothesis anticipates: recursive self-training on model-generated planning traces \emph{improves} capability rather than degrading it, provided a deterministic validator filters traces before reuse. This is not a contradiction so much as a bound: unfiltered synthetic contamination (this paper's setting) and validator-curated synthetic reuse (theirs) sit on opposite sides of exactly the strict-verifiability axis this paper's Claim 2 is built on --- their result is consistent with, and arguably strengthens, the correctness-strictness account, since it shows the same axis (mechanically checkable correctness) determines whether self-generated data helps or hurts, depending on whether verification is present. \citet{JointSelectionLarge} independently motivates why a single global filtering score is inadequate at trillion-token scale, showing single-metric (quality-only or diversity-only) data selection is fragile --- the same data-curation-economics stakes this paper's filtering-priority contribution addresses, but for a validated per-domain ranking rather than a single global score.

\section{Conclusion}

This work reframes a widely assumed governance simplification --- that synthetic-contamination damage is a property of the corpus as a whole --- as a testable claim about domain-selective fragility, rather than treating a single global mixing ratio as sufficient. The central methodological wager is that a cross-domain fragility ranking is only trustworthy once it is normalized against each domain's own incommensurable ground truth, isolated from incidental textual degeneracy, and shown to survive both a model-scale change and a contamination-generator swap; weaker designs risk reporting a ranking that is really an artifact of benchmark choice, text degeneracy, or generator idiosyncrasy. Whichever outcome the tested regime produces, the finding is decision-relevant: domain-selective damage would mean a uniform filtering policy systematically misallocates a finite curation budget, while a domain-uniform result would itself justify simplifying rather than complicating current practice. Two limitations bound this claim from the outset and are natural next steps rather than afterthoughts: the degeneracy-matched control rules out damage attributable to \emph{measured} textual degeneracy specifically, leaving latent, unmeasured quality defects as a residual, unresolved confound with syntheticness; and every claim here is scoped to the tested regime --- up to the low single-digit billions of parameters and the four specific domain corpora used --- so generalization to production-scale pretraining or to domains outside the tested four remains an open, explicitly flagged extrapolation rather than an implicit assumption.

\vspace{1em}
\hrule
\vspace{1em}

\section*{Pre-handoff checklist}
\begin{itemize}[leftmargin=1.4em]
\item \textbf{Organization/navigation} --- narrative flows problem $\to$ method $\to$ experiments $\to$ related work without forward references to undefined terms.
\item \textbf{Terminology discipline} --- ``correctness strictness,'' ``domain-selective,'' ``normalized capability delta,'' ``degeneracy-matched control,'' ``vulnerability ranking'' each defined at first use in Problem Definition or Methodology, used consistently thereafter; no internal stage IDs (\texttt{sp\_*}, \texttt{m\_*}, \texttt{rq\_*}, C1--C4) leaked into prose --- checked directly against the discussion pipeline log.
\item \textbf{Notation with worked example} --- the formal setup section defines $(D, \mathbb{D}, G, \kappa, \theta, S_d, \Delta_d, \beta_d)$ and the Problem Definition/Methodology sections consistently reuse this notation; no separate worked numeric example is included since no numbers exist yet (would require fabricating results).
\item \textbf{RQ framing matches outcomes} --- all 5 RQs are graded/quantitative (interaction significance, directional contrasts, rank-agreement statistics), each answered in its own subsection with an explicit falsifiable expected outcome and an explicit falsification condition, not a binary yes/no.
\item \textbf{No typesetting artifacts} --- only explicit placeholder markers in Results, per the drafting convention this is a pre-registered spec, not a report.
\item \textbf{Artifact plan} --- the contamination-injection protocol, degeneracy-logging instrumentation, and shard manifests are the natural release artifacts; not yet built, so not claimed as available.
\end{itemize}

\section*{Hand-off}
This is a draft, not a finished paper --- experiments have not run. Natural next steps: prose polish and venue-specific LaTeX formatting once a target venue is chosen, or revising an earlier stage if piloting surfaces a problem (e.g., if the walking skeleton surfaces that no domain shows a measurable delta at full contamination).

\bibliographystyle{plainnat}
\bibliography{refs}

\end{document}
